\documentclass[a4paper, 12pt]{article}

\usepackage{utils}

\renewcommand*{\today}{11 March 2026}

\begin{document}

\hotbox{Géometrie 2}{CM 4}{\today}

\section{Le triangle}

\subsection{Droites remarquables d'un triangle}

\begin{definition}
    Soit $ABC$ un triangle non plat. On appelle \textbf{cercle circonscrit} du triangle $ABC$ le cercle qui passe par les trois sommets $A$, $B$ et $C$.
\end{definition}

\begin{definition}
    Soit $ABC$ un triangle. On appelle \textbf{mediatrice} du segment $[BC]$ la droite qui est perpendiculaire à la droite $(BC)$ et qui passe par le milieu de $[BC]$.
\end{definition}

\begin{proposition}{}{}
    La mediatrice du segment $[BC]$ est l'ensemble des points équidistants de $B$ et de $C$.
\end{proposition}

\begin{theoreme}{}{}
    Les trois mediatrices d'un triangle sont concourantes. Le point de concours est appelé le \textbf{centre du cercle circonscrit} du triangle.
    
    Il est souvent noté $O$.
\end{theoreme}

\tikzexternalenable
\begin{figure}[H]
    \centering
    \resizebox{0.44\textwidth}{!}{%
    \begin{tikzpicture}
    \tikzset{
        mark 1/.style={postaction={decorate, decoration={markings, mark=at position 0.5 with {\draw[thick] (0,-3.5pt) -- (0,3.5pt);}}}},
        mark 2/.style={postaction={decorate, decoration={markings, mark=at position 0.5 with {\draw[thick] (-1.5pt,-3.5pt) -- (-1.5pt,3.5pt); \draw[thick] (1.5pt,-3.5pt) -- (1.5pt,3.5pt);}}}},
        mark 3/.style={postaction={decorate, decoration={markings, mark=at position 0.5 with {\draw[thick] (-2.5pt,-3.5pt) -- (-2.5pt,3.5pt); \draw[thick] (0,-3.5pt) -- (0,3.5pt); \draw[thick] (2.5pt,-3.5pt) -- (2.5pt,3.5pt);}}}}
    }

    \coordinate (O) at (0,0);
    \def\R{4}

    \coordinate (A) at (110:\R);
    \coordinate (B) at (210:\R);
    \coordinate (C) at (330:\R);

    \coordinate (Mc) at ($(A)!0.5!(B)$); % Milieu de AB
    \coordinate (Ma) at ($(B)!0.5!(C)$); % Milieu de BC
    \coordinate (Mb) at ($(A)!0.5!(C)$); % Milieu de AC

    \draw[red, thick] (O) circle (\R);

    \draw[blue, thin] ($(O)!-0.3!(Mc)$) -- ($(O)!1.3!(Mc)$);
    \draw[blue, thin] ($(O)!-0.3!(Ma)$) -- ($(O)!1.3!(Ma)$);
    \draw[blue, thin] ($(O)!-0.3!(Mb)$) -- ($(O)!1.3!(Mb)$);

    \draw[black, very thick] (A) -- (B) -- (C) -- cycle;

    \pic [draw, angle radius=3.5mm] {right angle = O--Mc--A};
    \pic [draw, angle radius=3.5mm] {right angle = O--Ma--C};
    \pic [draw, angle radius=3.5mm] {right angle = O--Mb--A};

    \draw[mark 1] (A) -- (Mc); \draw[mark 1] (Mc) -- (B);
    \draw[mark 2] (B) -- (Ma); \draw[mark 2] (Ma) -- (C);
    \draw[mark 3] (A) -- (Mb); \draw[mark 3] (Mb) -- (C);

    \foreach \p in {A,B,C,O}
        \filldraw[fill=white, draw=black, thick] (\p) circle (2.5pt);

    \node[above=3pt] at (A) {A};
    \node[below left=3pt] at (B) {B};
    \node[below right=3pt] at (C) {C};
    \node[below right=0pt, yshift=-2pt] at (O) {O};

\end{tikzpicture}
    }
    \caption{Cercle circonscrit (en rouge) et médiatrices (en bleu) d'un triangle $ABC$. Les trois médiatrices se coupent en un point $O$ qui est le centre du cercle circonscrit.}
\end{figure}
\tikzexternaldisable

\begin{definition}
    Soit $ABC$ un triangle non plat. On appelle \textbf{cercle inscrit} du triangle $ABC$ le cercle qui est tangent aux trois côtés du triangle.
\end{definition}

\begin{definition}
    Soit $ABC$ un triangle. On appelle \textbf{bissectrice} issue de $A$ la droite qui passe par $A$ et qui partage l'angle $\widehat{BAC}$ en deux angles de même mesure.
\end{definition}

\begin{proposition}{}{}
    La bissectrice issue de $A$ est l'ensemble des points équidistants des droites $(AB)$ et $(AC)$.
\end{proposition}

\begin{theoreme}{}{}
    Les trois bissectrices d'un triangle sont concourantes. Le point de concours est appelé le \textbf{centre du cercle inscrit} du triangle.
    
    Il est souvent noté $I$.
\end{theoreme}

\tikzexternalenable
\begin{figure}[H]
    \centering
    \resizebox{0.6\textwidth}{!}{%
    \begin{tikzpicture}

    \tikzset{
        mark segment/.style={postaction={decorate, decoration={markings, mark=at position 0.5 with {\draw[thick] (-1.5pt,-3.5pt) -- (-1.5pt,3.5pt); \draw[thick] (1.5pt,-3.5pt) -- (1.5pt,3.5pt);}}}},
        tick 1/.style={
            decoration={markings, mark=at position 0.5 with {\draw[thick] (0,-3.5pt) -- (0,3.5pt);}},
            postaction={decorate}
        },
        tick 2/.style={
            decoration={markings, mark=at position 0.5 with {\draw[thick] (-1.5pt,-3.5pt) -- (-1.5pt,3.5pt); \draw[thick] (1.5pt,-3.5pt) -- (1.5pt,3.5pt);}},
            postaction={decorate}
        },
        tick 3/.style={
            decoration={markings, mark=at position 0.5 with {\draw[thick] (-2.5pt,-3.5pt) -- (-2.5pt,3.5pt); \draw[thick] (0,-3.5pt) -- (0,3.5pt); \draw[thick] (2.5pt,-3.5pt) -- (2.5pt,3.5pt);}},
            postaction={decorate}
        }
    }

    % --- Définition des points principaux ---
    \coordinate (A) at (2.5, 6);
    \coordinate (B) at (0, 0);
    \coordinate (C) at (7, 0);
    
    % Centre du cercle inscrit (I) et rayon
    \coordinate (I) at (3, 2); 
    \def\r{2} 

    % --- Tracé du cercle inscrit ---
    \draw[red, thick] (I) circle (\r);

    % --- Calcul des intersections pour les bissectrices ---
    \path[name path=sidea] (B) -- (C);
    \path[name path=sideb] (A) -- (C);
    \path[name path=sidec] (A) -- (B);

    \path[name path=lineA] (A) -- ($(A)!2!(I)$);
    \path[name path=lineB] (B) -- ($(B)!2!(I)$);
    \path[name path=lineC] (C) -- ($(C)!2!(I)$);

    \path [name intersections={of=lineA and sidea, by=IntA}];
    \path [name intersections={of=lineB and sideb, by=IntB}];
    \path [name intersections={of=lineC and sidec, by=IntC}];

    % --- Construction de l'équidistance au sommet C ---
    % 1. On place un point P sur la bissectrice CI (à 60% de la distance)
    \coordinate (P) at ($(C)!0.4!(I)$);
    
    % 2. On crée une ligne perpendiculaire à CI passant par P
    \coordinate (Pdir1) at ($(P)!2cm!90:(I)$);
    \coordinate (Pdir2) at ($(P)!2cm!-90:(I)$);
    \path[name path=perpC] (Pdir1) -- (Pdir2);
    
    % 3. On trouve les intersections D et E avec les côtés AC et BC
    \path [name intersections={of=perpC and sideb, by=D}]; 
    \path [name intersections={of=perpC and sidea, by=E}]; 

    % --- Tracé du triangle et des bissectrices ---
    \draw[black, very thick] (A) -- (B) -- (C) -- cycle;

    \draw[blue, thin] (A) -- ($(A)!1.15!(IntA)$);
    \draw[blue, thin] (B) -- ($(B)!1.15!(IntB)$);
    \draw[blue, thin] (C) -- ($(C)!1.15!(IntC)$);

    % --- Dessin de la propriété au sommet C ---
    \draw[orange, thick] (D) -- (E); % Le segment perpendiculaire
    \pic [draw, angle radius=2.5mm] {right angle = I--P--D}; % L'angle droit
    \draw[mark segment] (D) -- (P); % Marque d'égalité PD
    \draw[mark segment] (P) -- (E); % Marque d'égalité PE

    % --- Nouvelles marques des angles (façon "arc barré") ---
    % Sommet B (1 trait)
    \pic [draw, angle radius=8mm] {angle = C--B--I};
    \pic [draw, angle radius=6mm] {angle = I--B--A};

    % Sommet C (2 traits)
    \pic [draw, angle radius=8mm] {angle = A--C--I};
    \pic [draw, angle radius=7mm] {angle = A--C--I};
    \pic [draw, angle radius=6mm] {angle = I--C--B};
    \pic [draw, angle radius=5mm] {angle = I--C--B};

    % Sommet A (3 traits)
    \pic [draw, angle radius=9mm] {angle = B--A--I};
    \pic [draw, angle radius=8mm] {angle = B--A--I};
    \pic [draw, angle radius=7mm] {angle = B--A--I};
    \pic [draw, angle radius=6mm] {angle = I--A--C};
    \pic [draw, angle radius=5mm] {angle = I--A--C};
    \pic [draw, angle radius=4mm] {angle = I--A--C};

    % --- Ajout des points et labels ---
    \foreach \p in {A,B,C,I}
        \filldraw[fill=white, draw=black, thick] (\p) circle (2pt);

    \node[above=3pt] at (A) {A};
    \node[below left=3pt] at (B) {B};
    \node[below right=3pt] at (C) {C};
    \node[above left=1pt, yshift=2pt] at (I) {I};


\end{tikzpicture}
    }
    \caption{Cercle inscrit (en rouge) et bissectrices (en bleu) d'un triangle $ABC$. Les trois bissectrices se coupent en un point $I$ qui est le centre du cercle inscrit.}
\end{figure}
\tikzexternaldisable

\begin{definition}
    Soit $ABC$ un triangle. On appelle \textbf{médiane} issue de $A$ la droite qui passe par $A$ et par le milieu de $[BC]$.
\end{definition}

\begin{proposition}{}{}
    Le centre de gravité d'un triangle est situé à $\frac{2}{3}$ de la médiane issue de chaque sommet à partir du sommet.
\end{proposition}

\begin{theoreme}{}{}
    Les trois médianes d'un triangle sont concourantes. Le point de concours est appelé le \textbf{centre de gravité} du triangle.
    
    Il est souvent noté $G$.
\end{theoreme}

\tikzexternalenable
\begin{figure}[H]
    \centering
    \resizebox{0.6\textwidth}{!}{%
    \begin{tikzpicture}

    % --- Styles robustes pour les marques de segments ---
    \tikzset{
        % Permet de dessiner 1, 2 ou 3 petits traits au milieu du segment
        mark 1/.style={postaction={decorate, decoration={markings, mark=at position 0.5 with {\draw[thick] (0,-3.5pt) -- (0,3.5pt);}}}},
        mark 2/.style={postaction={decorate, decoration={markings, mark=at position 0.5 with {\draw[thick] (-1.5pt,-3.5pt) -- (-1.5pt,3.5pt); \draw[thick] (1.5pt,-3.5pt) -- (1.5pt,3.5pt);}}}},
        mark 3/.style={postaction={decorate, decoration={markings, mark=at position 0.5 with {\draw[thick] (-2.5pt,-3.5pt) -- (-2.5pt,3.5pt); \draw[thick] (0,-3.5pt) -- (0,3.5pt); \draw[thick] (2.5pt,-3.5pt) -- (2.5pt,3.5pt);}}}}
    }

    % --- Définition des sommets du triangle ---
    \coordinate (A) at (1.5, 5);
    \coordinate (B) at (0, 0);
    \coordinate (C) at (7, 0);

    % --- Calcul des milieux des côtés ---
    \coordinate (Aprime) at ($(B)!0.5!(C)$); % Milieu de BC (opposé à A)
    \coordinate (Bprime) at ($(A)!0.5!(C)$); % Milieu de AC (opposé à B)
    \coordinate (Cprime) at ($(A)!0.5!(B)$); % Milieu de AB (opposé à C)

    % --- Calcul du Centre de Gravité (G) ---
    % TikZ sait calculer les barycentres nativement, c'est parfait pour G !
    \coordinate (G) at (barycentric cs:A=1,B=1,C=1);

    % --- Tracé du triangle ---
    \draw[black, very thick] (A) -- (B) -- (C) -- cycle;

    % --- Tracé des médianes ---
    \draw[blue, thin] (A) -- (Aprime);
    \draw[blue, thin] (B) -- (Bprime);
    \draw[blue, thin] (C) -- (Cprime);

    % --- Ajout des marques de segments (équidistance) ---
    % Côté BC (coupé par A') -> 1 trait
    \draw[mark 1] (B) -- (Aprime); 
    \draw[mark 1] (Aprime) -- (C);
    
    % Côté AC (coupé par B') -> 2 traits
    \draw[mark 2] (A) -- (Bprime); 
    \draw[mark 2] (Bprime) -- (C);

    % Côté AB (coupé par C') -> 3 traits
    \draw[mark 3] (A) -- (Cprime); 
    \draw[mark 3] (Cprime) -- (B);

    % --- Ajout des points et labels ---
    \foreach \p in {A,B,C,G,Aprime,Bprime,Cprime}
        \filldraw[fill=white, draw=black, thick] (\p) circle (2pt);

    \node[above=3pt] at (A) {A};
    \node[below left=3pt] at (B) {B};
    \node[below right=3pt] at (C) {C};
    
    % Labels des milieux
    \node[below=3pt] at (Aprime) {A'};
    \node[above right=2pt] at (Bprime) {B'};
    \node[above left=2pt] at (Cprime) {C'};

    % Label du centre de gravité
    \node[above right=0pt, yshift=2pt] at (G) {G};

\end{tikzpicture}
    }
    \caption{Médianes (en bleu) d'un triangle $ABC$. Les trois médianes se coupent en un point $G$ qui est le centre de gravité du triangle.}
\end{figure}
\tikzexternaldisable

\begin{definition}
    Soit $ABC$ un triangle. On appelle \textbf{hauteur} issue de $A$ la droite qui passe par $A$ et qui est perpendiculaire à la droite $(BC)$.
\end{definition}

\begin{theoreme}{}{}
    Les trois hauteurs d'un triangle sont concourantes. Le point de concours est appelé l'\textbf{orthocentre} du triangle.
    
    Il est souvent noté $H$.
\end{theoreme}

\tikzexternalenable
\begin{figure}[H]
    \centering
    \resizebox{0.6\textwidth}{!}{%
    \begin{tikzpicture}

    % --- Définition des sommets du triangle ---
    % On choisit un triangle acutangle (tous les angles < 90°) 
    % pour que l'orthocentre soit bien à l'intérieur de la figure.
    \coordinate (A) at (2, 5);
    \coordinate (B) at (0, 0);
    \coordinate (C) at (7, 0);

    % --- Calcul des pieds des hauteurs par projection orthogonale ---
    % La syntaxe magique de TikZ : $(B)!(A)!(C)$ projette le point A sur la droite (BC)
    \coordinate (Ha) at ($(B)!(A)!(C)$); 
    \coordinate (Hb) at ($(A)!(B)!(C)$); 
    \coordinate (Hc) at ($(A)!(C)!(B)$);

    % --- Calcul de l'Orthocentre (H) ---
    % On le trouve en cherchant l'intersection de deux hauteurs
    \path[name path=altA] (A) -- (Ha);
    \path[name path=altB] (B) -- (Hb);
    \path [name intersections={of=altA and altB, by=H}];

    % --- Tracé du triangle ---
    \draw[black, very thick] (A) -- (B) -- (C) -- cycle;

    % --- Tracé des hauteurs ---
    \draw[blue, thin] (A) -- (Ha);
    \draw[blue, thin] (B) -- (Hb);
    \draw[blue, thin] (C) -- (Hc);

    % --- Ajout des symboles d'angles droits ---
    \pic [draw, angle radius=3mm] {right angle = A--Ha--C};
    \pic [draw, angle radius=3mm] {right angle = B--Hb--C};
    \pic [draw, angle radius=3mm] {right angle = C--Hc--B};

    % --- Ajout des points et labels ---
    \foreach \p in {A,B,C,H}
        \filldraw[fill=white, draw=black, thick] (\p) circle (2pt);

    \node[above=3pt] at (A) {A};
    \node[below left=3pt] at (B) {B};
    \node[below right=3pt] at (C) {C};
    
    % Labels des pieds des hauteurs
    \node[below=3pt] at (Ha) {$H_A$};
    \node[above right=1pt] at (Hb) {$H_B$};
    \node[above left=1pt] at (Hc) {$H_C$};

    % Label de l'orthocentre
    \node[above left=1pt, yshift=2pt] at (H) {H};

\end{tikzpicture}
    }
    \caption{Hauteurs (en bleu) d'un triangle $ABC$. Les trois hauteurs se coupent en un point $H$ qui est l'orthocentre du triangle.}
\end{figure}
\tikzexternaldisable

\subsection{Les propriétés élémentaires}

\end{document}